% GENERATED from paper/draft/00_front_matter.md — do not edit
Organisations must decide which decisions an AI agent may take without
human review. Correct outcomes alone do not show whether an agent
responds appropriately to evidence. We define decision-structured agency
(DSA) as responsiveness to the evidence that should determine a class of
decisions, measured by removing documents, reversing their claims or
adding independent evidence. In a pre-registered pilot on synthetic
counterparty verification (14 model configurations, eight developers),
we compare delegation on DSA, track records and an LLM judge's reading
of the decision trace. Deterministic programs implementing our reading
of the policy generate the answer key. The registered test favoured DSA
over track records at matched coverage (one-sided p \textless{} 0.001;
configuration bootstrap p = 0.002). The advantage, however, depended on
one intervention: removing a register excerpt the policy did not
require. Excluding two interventions flagged by a policy audit (p =
0.20) or revising the key after the author's partial review (p = 0.28)
removed statistical support for the advantage. Across 60 combinations of
policy readings, it survived only under readings stricter than the
underlying legal texts. Under every key, the judge passed most decisions
that were correct on the original file but failed an evidence test
(59--70\%); a stronger judge did no better. Much of the judge-based
rules' deficit came from one class where the main judge read the policy
differently from the programs. The primary result depended on how the
policy was translated into an answer key. That key is part of the
instrument and requires independent scrutiny.
